\bigskip\Needspace{10\baselineskip}
\pdfbookmark[1]{Pallas Olbersiana. Dezember 1802}{pallas180212}
\begin{center}
PALLAS OLBERSIANA.\par
\medskip\hrule\medskip
Monatliche Correspondenz. Band VI. S.\ 579—583. Dezember 1802.\par
\medskip\hrule
\end{center}\medskip
—\quad —\quad —

Dr.\ \textsc{Gauß} hat nun alle diese Mailänder Beobachtungen \textsc{Oriani}s mit seinen Elementen III für die Bahn dieses Planeten [S.\ 218 dieses Bandes] verglichen und die folgende Übereinstimmung gefunden.
\Needspace{22\baselineskip}
\begin{center}\textit{Vergleich der Beobachtungen \textsc{Oriani}s am Mailänder Äquatorialsektor mit den Bahnelementen III von Dr.\ \textsc{Gauß}.}\end{center}
\begin{center}
\begingroup\footnotesize\setlength{\tabcolsep}{2pt}
\begin{tabular}{l|r|r|r|r|r}
1802&\shortstack{Mittl. Zeit\\in Mailand}&\shortstack{Berechnete\\RA der Pallas}&\shortstack{Unterschied\\R—B}&\shortstack{Berechnete nördl.\\Dekl. der $\PallasSymbol$}&\shortstack{Unterschied\\R—B}\\\hline
Juli 10&$9^{\mathrm u}\ 40^{\mathrm m}\ 23^{\mathrm s}$&$191^0\ 27'\ 33''$&$\left\{\begin{array}{r}-11''\\-35\end{array}\right.$&$17^0\ 56'\ 20''$&$\left\{\begin{array}{r}-4''\\+4\end{array}\right.$\\
18&$9\ 36\ 4$&$193\ 43\ 47$&$-17$&$17\ 2\ 29$&$-38$\\
24&$9\ 32\ 16$&$195\ 30\ 59$&$-29$&$16\ 19\ 59$&$+12$\\
25&$9\ 44\ 12$&$195\ 49\ 25$&$-23$&$16\ 12\ 40$&$+12$\\
28&$9\ 8\ 6$&$196\ 44\ 25$&$\left\{\begin{array}{r}-33\\-28\end{array}\right.$&$15\ 50\ 56$&$\left\{\begin{array}{r}+1\\+16\end{array}\right.$\\
29&$9\ 11\ 28$&$197\ 3\ 9$&$\left\{\begin{array}{r}-36\\-22\\-36\end{array}\right.$&$15\ 43\ 33$&$\left\{\begin{array}{r}+1\\+18\\+14\end{array}\right.$\\
31&$9\ 6\ 9$&$197\ 40\ 46$&$-34$&$15\ 28\ 55$&$-1$\\
Aug.\ 1&$9\ 4\ 14$&$197\ 59\ 44$&$\left\{\begin{array}{r}-40\\-45\\-49\end{array}\right.$&$15\ 21\ 18$&$\left\{\begin{array}{r}-20\\+2\\+5\end{array}\right.$\\
2&$9\ 23\ 37$&$198\ 19\ 4$&$\left\{\begin{array}{r}-36\\-41\end{array}\right.$&$15\ 13\ 43$&$\left\{\begin{array}{r}+10\\+12\end{array}\right.$
\end{tabular}\endgroup
\end{center}
\SourceLeaf{238}%
\begin{center}
\begingroup\footnotesize\setlength{\tabcolsep}{2pt}
\begin{tabular}{l|r|r|r|r|r}
1802&\shortstack{Mittl. Zeit\\in Mailand}&\shortstack{Berechnete\\RA der Pallas}&\shortstack{Unterschied\\R—B}&\shortstack{Berechnete nördl.\\Dekl. der $\PallasSymbol$}&\shortstack{Unterschied\\R—B}\\\hline
Aug.\ 4&$8^{\mathrm u}\ 55^{\mathrm m}\ 38^{\mathrm s}$&$198\ 57\ 9$&$\left\{\begin{array}{r}-39''\\-47\end{array}\right.$&$14\ 58\ 52$&$\left\{\begin{array}{r}-5''\\-4\end{array}\right.$\\
5&$9\ 14\ 57$&$199\ 16\ 47$&$-39$&$14\ 51\ 13$&$\left\{\begin{array}{r}+15\\+28\end{array}\right.$\\
6&$9\ 6\ 27$&$199\ 36\ 7$&$\left\{\begin{array}{r}-50\\-51\end{array}\right.$&$14\ 43\ 43$&$-3$\\
7&$9\ 4\ 18$&$199\ 55\ 38$&$\left\{\begin{array}{r}-48\\-47\end{array}\right.$&$14\ 36\ 8$&$+15$\\
8&$9\ 0\ 14$&$200\ 15\ 9$&$-56$&$14\ 28\ 35$&$+14$
\end{tabular}\endgroup
\end{center}

Dr.\ \textsc{Gauß} hat zwar später zum vierten Mal verbesserte Elemente angegeben (Oktoberheft, S.\ 394) [S.\ 223 dieses Bandes], sie aber nicht selbst mit den obigen Mailänder Beobachtungen verglichen. Er hat jedoch den folgenden Vergleich angestellt, anhand dessen man ungefähr beurteilen kann, inwieweit diese Elemente IV besser mit den Beobachtungen übereinstimmen als die Elemente III.
\begin{center}
Die Elemente IV ergeben\par\medskip
\begin{tabular}{l|r|r}
1802&Eine größere RA&Eine kleinere Deklination\\\hline
Juli 8&$24''$&$8''$\\
Aug.\ 4&$66$&$0$
\end{tabular}\par\medskip
als die Elemente III.
\end{center}

Gegen Ende der Beobachtungen liefern die Elemente III also eine um etwa 1 Minute zu kleine Rektaszension und die Elemente IV eine um etwa $20''$ zu große. Die Deklinationen stimmen bei beiden ziemlich gut überein, soweit es die Genauigkeit der Beobachtungen zulässt.

Anhand dieser jüngsten Beobachtungen \textsc{Oriani}s hat Dr.\ \textsc{Gauß} die Elemente dieser Bahn nun auch zum fünften Mal verbessert und die folgenden Werte erhalten.
\Needspace{14\baselineskip}
\begin{center}\textit{Elemente V der Pallas Olbersiana.}\end{center}
\begin{center}
\begingroup\small\setlength{\tabcolsep}{3pt}
\begin{tabular}{l r}
\shortstack[l]{Epoche 31. März 1802\\mittags in Seeberg}&$162^0\ 46'\ 58''2$\\
Aphel $\ .\ .\ .\ .\ .\ .\ .\ .\ .$&$301\ 28\ 24.0$\\
Aufsteigender Knoten $\ .\ .\ .\ .\ .\ .$&$172\ 27\ 3.0$\\
Mittlere tägliche tropische Bewegung $\ .$&$769''583$\\[3pt]
Logarithmus der großen Halbachse&$0.4425529$\\
Exzentrizität $\ .\ .\ .\ .\ .\ .\ .\ .\ .$&$0.244976$\\
Neigung der Bahn $\ .\ .\ .\ .\ .\ .\ .$&$34^0\ 37'\ 40''$
\end{tabular}\endgroup
\end{center}

Dr.\ \textsc{Gauß} verglich nun die späteste aller Beobachtungen \textsc{Messier}s mit diesen Elementen (Novemberheft, S.\ 502). Er fand die folgenden Werte.
\Needspace{12\baselineskip}
\begin{center}
\begingroup\small\setlength{\tabcolsep}{4pt}
\begin{tabular}{l|r|r|r}
1802&\shortstack{Mittlere Zeit\\in Paris}&\shortstack{Berechnete RA\\der Pallas}&\shortstack{Unterschied\\R—B}\\\hline
Sept.\ 21&$7^{\mathrm u}\ 28^{\mathrm m}$&$215^0\ 49'\ 14''$&$+28''$
\end{tabular}\par\medskip
\begin{tabular}{r|r|r|r}
\shortstack{Berechnete Deklination\\der Pallas}&\shortstack{Unterschied\\R—B}&\shortstack{Abstand\\von der $\EarthSymbol$}&Lichtstärke\\\hline
$9^0\ 0'\ 16''\ \mathrm N.$&$+48''$&$3.518$&$0.485$
\end{tabular}\endgroup
\end{center}

Der Unterschied zwischen Rechnung und Beobachtung ist also noch mäßig. Dr.\ \textsc{Gauß} schreibt daher Folgendes.
„Da diese Beobachtung \textsc{Messier}s bei so ungünstiger Lage gemacht wurde (nach einer überschlägigen Rechnung stand Pallas nur 13 Grad über dem Horizont), bin ich ganz Ihrer Meinung, dass man aus dieser \textit{einzelnen} Beobachtung wenig Verlässliches zur Verbesserung der Elemente gewinnen kann. Die Verbesserung dürfte jedoch auf die Position der \textit{Pallas} im kommenden Jahr keinen sehr großen Einfluss haben. Ich habe mehrere Positionen nach den Elementen V neu
\SourceLeaf{239}%
berechnet. Für den 4. Febr., 24. März, 11. Mai und 28. Juni 1803 fand ich die Rektaszensionen um $3\frac12$ Min., $5\frac12$ Min., $7\frac12$ Min. bzw. 9 Min. kleiner und die Deklinationen alle um ungefähr 1 Min. größer als in der nach den Elementen IV berechneten Ephemeride.“ (Oktoberheft, S.\ 395.) [S.\ 224 dieses Bandes.]

Um auch hier keinen Wunsch offenzulassen, hat Dr.\ \textsc{Gauß} die ganze Ephemeride für den Lauf der Pallas im Jahr 1803 nach diesen Elementen V neu berechnet und um einen weiteren Monat verlängert. Am 3. April 1803 wird \textit{Pallas} heller sein, als sie am 21. September dieses Jahres war. Da sie dann auch in viel größerer Höhe über dem Horizont beobachtet werden kann, ist jetzt wohl kaum noch daran zu zweifeln, dass wir sie mit unseren stärkeren Fernrohren sehen werden. Es folgt die Ephemeride. [Die Ephemeride vom 4. Febr.\ bis 9. Aug.\ 1803 mit Rektaszension, Deklination und Lichtstärke in der bisher verwendeten Form.]

In der dritten Spalte dient die Lichtstärke, die der Planet in der Entfernung 1 von Sonne und Erde hätte, als Einheit. Nach diesem Maßstab hatte die Lichtstärke im Jahr 1802 die folgenden Werte.
\begin{center}
\begingroup\small
\begin{tabular}{llr r l}
1802&April&4 $\ .\ .$&$0.08997$&[Handschrift.] Genauer $0.09019$\\
&Mai&16 $\ .\ .$&$0.04740$&\\
&Aug.&10 $\ .\ .$&$0.01455$&\\
&Sept.&21 $\ .\ .$&$0.01030$&
\end{tabular}\par\medskip
\begin{tabular}{llr r}
1801&Jan.&1 $\ .\ .$&$0.03637$\\
&Febr.&11 $\ .\ .$&$0.02409$\\
&Dec.&7 $\ .\ .$&$0.02494$\\
1802&Febr.&26 $\ .\ .$&$0.05633$
\end{tabular}\endgroup
\end{center}
\begin{flushright}[\textsc{von Zach.}]\end{flushright}

